% !TeX root = ../main.tex
\chapter[Introduction]{Introduction}
\label{chap:introduction}

\epigraph{\textit{I felt the weight of unmapped worlds, unborn language.}}{--- William Finnegan \\\footnotesize\textit{Barbarian Days: A Surfing Life}}

The Graphics Processing Unit, or GPU, has transcended its name and early use for only rendering graphics, as modern applications have harnessed its parallel processing ability to accelerate an ever-increasing number of general-purpose compute workloads (GPGPU)~\cite{owens2007gpgpu}. Their use in training convolutional neural networks helped trigger the deep learning revolution~\cite{krizhevsky2017imagenet}, and GPUs are currently used by a variety of machine learning applications~\cite{reuther2019survey}. GPUs help power fundamental applications like encryption~\cite{ozerk2022encrypt}, as well as transformational technologies like particle simulation for drug discovery~\cite{pandey2022transformational} and self-driving cars~\cite{bojarski2016endendlearningselfdriving}.

CUDA~\cite{cuda}, released by NVIDIA in 2007 and targeting only NVIDIA GPUs, was the first widely successful GPGPU framework, and highlighted the advantages of using a compute-focused language to accelerate general-purpose applications on GPUs. However, modern GPU computing spans devices and software stacks from multiple vendors, with hardware characteristics that vary across vendors, targeted platform (e.g., datacenter vs. desktop vs. mobile) and architecture generations within each vendor. To support cross-platform development and portable applications, GPU programming frameworks have been developed to provide common abstractions across devices. OpenCL established an early cross-vendor model, but prior empirical studies found portability issues that stopped applications from successfully executing on GPUs from different vendors~\cite{sorensen2016hitchhiker}. More recently, frameworks like Vulkan~\cite{vulkan} and WebGPU~\cite{webgpu} have established themselves as prominent choices for cross-platform GPU compatibility, while WebGPU additionally allows GPU compute applications to be developed for the browser. These abstractions now support emerging and complex workloads like cross-platform and browser-based large language model (LLM) inference libraries~\cite{ruan2024webllmhighperformanceinbrowserllm,transformersjsDocs,levine2026llamaweb}.

One of the fundamental aspects of performance-oriented Von Neumann style programming frameworks, e.g., those used for modern CPUs and GPUs, is how they define operations on memory. In programming frameworks where multiple operations run in parallel, e.g., on GPUs, these definitions must also specify the semantics when multiple processing elements (commonly called \textit{threads}) concurrently access the same underlying memory. This set of definitions and rules is called a \textit{memory consistency specification} (MCS). An MCS plays two related roles: (1) it is a contract that implementations must satisfy, i.e., an implementation is \textit{sound} with respect to an MCS if it never exhibits behaviors that the MCS forbids, and (2) it is an abstraction of real machines, i.e., an implementation may be sound while still being more restrictive than the MCS, since it need not exhibit every behavior that the MCS permits.

Modern processors include many features that complicate the development of MCSs, including instruction reordering, buffering, and deep cache hierarchies, causing memory operations to not necessarily become visible to other threads in intuitive orders~\cite{nagarajan2020primer,gharachorloo1991performance,sarkar2011understanding,sarkar2009semantics}. So, while early MCSs were described in informal natural language, CPU architectures like x86~\cite{owens2009x86} and ARM~\cite{alglave2009semantics} formalize their MCSs to more precisely define allowed behaviors. Languages like C++~\cite{batty2011mathematizing} and Java~\cite{manson2005java} have also formalized MCSs that abstract over the underlying hardware.

In comparison, cross-platform GPU frameworks have been developed more recently and target a diverse set of underlying architectures from many vendors with corner-case behaviors that are not well understood. While previous work formalized the MCS of virtual ISAs from specific GPU vendors like NVIDIA~\cite{lustig2019ptx} and developed a formal MCS for cross-platform frameworks like OpenCL~\cite{batty2016overhauling}, the empirically observable behaviors of GPUs from different vendors have been shown to vary substantially and numerous bugs in GPU software and hardware stacks have been found throughout the years~\cite{alglave2015gpu,alglave2010fences,kirkham2020foundations}. The bugs show that implementations require more rigorous testing techniques to validate soundness, while the gap between specified and observed behaviors can motivate more precise models that squeeze more performance out of existing GPUs.

Beyond ensuring that hardware and compilers conform to a given MCS, empirically testing and understanding the concurrent shared-memory behavior of cross-platform frameworks on a variety of GPUs is also important for application correctness and security. Early work on accelerating GPU applications simply assumed that non-intuitive MCS behaviors could not occur~\cite{xiao2010interblock}, while further work showed that these assumptions were flawed and could lead to bugs depending on the environment an application runs in~\cite{sorensen2016exposing}. The behaviors described by an MCS are often due to underlying hardware features, which also motivates the question of whether observations of these behaviors can reveal information about architectural state and application characteristics. For example, previous work has identified hardware side-channels due to memory and cache behaviors~\cite{jiang2024sync,guo2024usenix,wang2024gpuzip}, but no work has explicitly considered the MCS as a source of side-channel vulnerabilities.

An MCS is also responsible for defining which memory patterns are meaningful and which lead to harmful data races that fall outside the model, helping to determine where implementations and runtimes must be hardened to avoid memory safety vulnerabilities. Numerous exploits targeting various levels of GPU memory hierarchies have been found in recent years~\cite{sorensen2024leftoverlocals,pustelnik2024whispering,dipietro2016cudaleaks,guo2024usenix}, and security concerns are ratcheted up another level in the browser, where users often execute code from untrusted sources, and where GPU vulnerabilities have been used to steal webpages~\cite{lee2014stealing} and create side-channels~\cite{giner2024cacheattack}.

To understand existing techniques and insufficiencies in testing and validating an MCS, we now turn to a quick primer on memory consistency.

\newcommand{\mprlxintrolitmustest}{
    \begin{tabular}{l l}
    \midrule
    \multicolumn{1}{c}{\underline{thread 0}} & \multicolumn{1}{c}{\underline{thread 1}} \\
        \introMpEventA\hspace{.1cm} \texttt{store(data, 1)} &  \introMpEventC\hspace{.1cm} \texttt{let r0 = load(flag)}  \\
        \introMpEventB\hspace{.1cm} \texttt{store(flag, 1)} &
        \introMpEventD\hspace{.1cm} \texttt{let r1 = load(data)}  \\
        \midrule
        %\multicolumn{2}{c}{Condition: \texttt{r0 == 1 \&\& r1 == 0}}
    \end{tabular}
}

\newcommand{\introMpEventA}{\circledBlack{\vphantom{A}a}\xspace}
\newcommand{\introMpEventB}{\circledBlack{\vphantom{A}b}\xspace}
\newcommand{\introMpEventC}{\circledWhite{\vphantom{A}c}\xspace}
\newcommand{\introMpEventD}{\circledWhite{\vphantom{A}d}\xspace}
\newcommand{\introMpFenceX}{\circledBlack{\vphantom{A}x}\xspace}
\newcommand{\introMpFenceY}{\circledWhite{\vphantom{A}y}\xspace}

\newcommand{\mprelacqintrolitmustest}{
    \begin{tabular}{l l}
    \midrule
    \multicolumn{1}{c}{\underline{thread 0}} & \multicolumn{1}{c}{\underline{thread 1}} \\
        \introMpEventA\hspace{.1cm} \texttt{store(data, 1)} &  \introMpEventC\hspace{.1cm} \texttt{let r0 = load(flag)}  \\
        \introMpFenceX\hspace{.1cm} \texttt{fence(release)} &  \introMpFenceY\hspace{.1cm} \texttt{fence(acquire)}  \\
        \introMpEventB\hspace{.1cm} \texttt{store(flag, 1)} &
        \introMpEventD\hspace{.1cm} \texttt{let r1 = load(data)}  \\
        \midrule
        %\multicolumn{2}{c}{Condition: \texttt{r0 == 1 \&\& r1 == 0}}
    \end{tabular}
}


\begin{figure}
\centering
\small
\captionsetup[subfigure]{width=.7\columnwidth}
\subfloat[Message passing test.]{\resizebox{.7\columnwidth}{!}{\mprlxintrolitmustest} \label{fig:intro-mprlx}} \\
\subfloat[Message passing test with fences.]{\resizebox{.7\columnwidth}{!}{\mprelacqintrolitmustest} \label{fig:intro-mprelacq}}
\caption{Example of a \textit{message passing} litmus test.}
\label{fig:intro-mp}
\end{figure}



\section{A Primer on Memory Consistency}

Ever since computer scientists had the temerity to combine multiple processors that shared access to a common memory region with private caches for each processor, memory consistency has vexed and confounded hardware and language designers alike. In fact, memory consistency has been called the rocket science of computer science by none other than the designer of Linux~\cite{torvalds2021volatileif}. Yet, perhaps not unlike modern Don Quixotes tilting at windmills, computer scientists continue charging headlong into the fray of building and testing models of memory consistency, and it is this quest that we now join as well.

\myfiglong~\ref{fig:intro-mprlx} shows a small, two-thread program, called a \textit{litmus test} in memory consistency literature, that operates on memory locations \texttt{data} and \texttt{flag}. Assume for now that these operations are direct hardware load and store instructions on memory locations, executed on a processor with a cache hierarchy that includes private caches for each thread. This program is a form of \textit{message passing}; the first thread writes some value to \texttt{data} (in this case 1), then signals that \texttt{data} is ready by writing the value 1 to \texttt{flag}. In parallel, the second thread first reads \texttt{flag} and then reads \texttt{data}.

The most intuitive way to analyze this program is to consider the possible orderings while respecting the fact that event \introMpEventA occurs before \introMpEventB on thread 0 and event \introMpEventC occurs before \introMpEventD on thread 1. This model of execution is known as \textit{sequential consistency} (SC)~\cite{lamport1979multiprocessor}. The first six rows of \mytab\ref{tab:intro-mp-orderings} show possible outcomes under this model. The first two rows show schedules where each thread executes sequentially; if thread 0 executes first, then both \texttt{r0} and \texttt{r1} end up with the value 1, whereas if the orders are reversed, both end up with the value 0. Given that the threads are operating in parallel, it is also possible that the events are interleaved. The middle 4 rows show possible interleaved schedules; in all cases, \texttt{r0} ends up with the value 0 and \texttt{r1} ends up with the value 1.

\begin{table}[t]
\centering
\begin{tabular}{l l l}
\toprule
\multicolumn{1}{c}{\textbf{Event Ordering}} & \multicolumn{1}{c}{\textbf{Outcome}} & \multicolumn{1}{c}{\textbf{Behavior}} \\
\midrule
\introMpEventA $\rightarrow$ \introMpEventB $\rightarrow$ \introMpEventC $\rightarrow$ \introMpEventD & \texttt{r0 == 1 \&\& r1 == 1} & \multirow{2}{*}{\texttt{Sequential}} \\
\introMpEventC $\rightarrow$ \introMpEventD $\rightarrow$ \introMpEventA $\rightarrow$ \introMpEventB & \texttt{r0 == 0 \&\& r1 == 0} & \\
\midrule
\introMpEventA $\rightarrow$ \introMpEventC $\rightarrow$ \introMpEventB $\rightarrow$ \introMpEventD & \multirow{4}{*}{\texttt{r0 == 0 \&\& r1 == 1}} & \multirow{4}{*}{\texttt{Interleaved}} \\
\introMpEventA $\rightarrow$ \introMpEventC $\rightarrow$ \introMpEventD $\rightarrow$ \introMpEventB & & \\
\introMpEventC $\rightarrow$ \introMpEventA $\rightarrow$ \introMpEventD $\rightarrow$ \introMpEventB & & \\
\introMpEventC $\rightarrow$ \introMpEventA $\rightarrow$ \introMpEventB $\rightarrow$ \introMpEventD & & \\
\midrule
\introMpEventB $\rightarrow$ \introMpEventC $\rightarrow$ \introMpEventD $\rightarrow$ \introMpEventA & \multirow{2}{*}{\texttt{r0 == 1 \&\& r1 == 0}} & \multirow{2}{*}{\texttt{Weak}} \\
\introMpEventD $\rightarrow$ \introMpEventA $\rightarrow$ \introMpEventB $\rightarrow$ \introMpEventC & & \\
\bottomrule
\end{tabular}
\caption{Example orderings for the message-passing events in \myfig\ref{fig:intro-mprlx}.}
\label{tab:intro-mp-orderings}
\end{table}

Unfortunately, almost all modern multiprocessors, including GPUs, are not sequentially consistent. Due to out-of-order processing, private caches, load and store buffers, and other proprietary features that increase overall performance~\cite{gharachorloo1991performance}, the order in which threads observe the effects of other threads operating on memory is non-deterministic and governed by many complex rules. A non-SC MCS is called a \textit{relaxed}~\cite{adve1996relaxed} or \textit{weak}~\cite{dubois1986weak} model, today often interchangeably. For consistency (no pun intended) in this thesis a non-SC model is referred to as \textit{relaxed}, and the non-SC behaviors allowed by that model are referred to as \textit{weak behaviors}.

The bottom two rows of \mytab\ref{tab:intro-mp-orderings} show the weak behaviors, where \texttt{r0} ends up with the value 1 but \texttt{r1} ends up with the value 0. Note that for this outcome to occur, it must appear as if either \introMpEventA is reordered with respect to \introMpEventB, or \introMpEventC is reordered with respect to \introMpEventD, e.g., due to the order in which cache lines are flushed to main memory. Given a set of plain reads and writes on hardware, the results in this thesis show that modern GPUs from every major vendor will allow this outcome to occur (\mychap\ref{chap:gpuharbor}).

In many scenarios, the weak outcome is not desirable. Suppose that a programmer would like to use \texttt{flag} to signal to another thread that some value is ready; given the possibility of a weak outcome using plain reads and writes, there is no way for the thread to know whether the value of \texttt{data} is valid. Therefore, modern hardware also implements special instructions called \textit{fences}, which signal that memory operations should be propagated in certain orders. For example, \myfig\ref{fig:intro-mprelacq} shows the same message passing program, but this time with fences inserted between the events on each thread. Specifically, the fence on thread 0 is a \textit{release} fence, which indicates to the hardware that any reads or writes that occur before the write event \introMpEventB must appear to occur before \introMpEventB itself, from the viewpoint of other threads. The fence on thread 1 is a complementary \textit{acquire} fence, which indicates that any read or write that occurs after the read event \introMpEventC must appear to occur after \introMpEventC. Under the constraints added by these fences, it is no longer possible for the weak behavior to legally occur, giving programmers the ability to reliably transfer data across threads.

\paragraph{Higher-Level MCSs.}
So far, the assumption has been that the store and load instructions in the message passing program have been occurring directly on hardware. Of course, most code is not written to run directly on hardware, but rather in higher-level representations that are then compiled into formats that hardware can run. Languages like C++ and Java have formalized their own MCSs~\cite{batty2011mathematizing,manson2005java} and shown that they are compatible with MCSs for CPUs like x86~\cite{owens2009x86}, Arm~\cite{alglave2009semantics}, and Power~\cite{alglave2009semantics}. In the GPU ecosystem, NVIDIA has formalized an MCS for their higher-level PTX instruction set~\cite{lustig2019ptx}, but other GPU vendors lag behind on formalizing MCSs for their specific software and hardware stacks. This has led cross-platform GPU frameworks like Vulkan and WebGPU to provide a lowest-common-denominator MCS and depend on informal specifications and committee discussions.

The presence of optimizing compilers between high-level programming languages and the code that actually executes on hardware complicates the development of the language-level MCS. Compilers look for opportunities to eliminate and reorder instructions while maintaining the semantics of the original program. Memory reads and writes are among the expensive operations compilers seek to optimize, for example by removing redundant load or store instructions or reordering those instructions to take better advantage of hardware pipelining. However, these optimizations are often in conflict with the goals of programmers seeking to explicitly communicate across threads using memory operations. Therefore, modern languages like C++ (along with many GPU languages) have introduced \textit{atomic} memory operations, which signal to the compiler that these operations' associated memory locations may be accessed concurrently by multiple threads. For the example in \myfig\ref{fig:intro-mprelacq}, each read and write would need to be an atomic operation in a high-level language for the message passing program to be well-formed.

If the instructions in \myfig\ref{fig:intro-mp} are written in a high-level language but not specified to be atomic, then a \textit{data race} occurs. The effects of a data race in a high-level language include torn reads or writes, impossible values occurring, and infinite loops/deadlocks~\cite{boehm2011benign,dolan2018bounding}, leading to languages like C++ giving ``catch-fire'', or undefined semantics to data races. Undefined behavior is a natural source of potential security vulnerabilities, as implementations may not have protections against all the possible (and unknown) effects. Therefore, languages like Java give weak semantics to data races, although problems in the Java memory model continue to be discovered~\cite{panneke2026jmt}. Browser-based CPU languages like JavaScript and WebAssembly also give very weak semantics to data races~\cite{watt2020repairing,watt2019weakening}, as security in the browser where users regularly execute untrusted code is extremely critical. However, GPU languages have not generally grappled with the complexity of the security concerns with data races, which is especially concerning given the large number of GPU memory safety vulnerabilities found in recent years~\cite{lee2014stealing,donaldson2017graphicsfuzz,sorensen2024leftoverlocals,pustelnik2024whispering,guo2024usenix}.

\section{Validation Challenges}

Analyzing small programs like the ones in \myfig\ref{fig:intro-mp} is useful for understanding what a given MCS does and does not allow, but it does not give any insight into what actually occurs when the program is run on real hardware. As GPU usage continues to increase, and cross-platform GPU frameworks gain more prominence, ensuring application correctness and safety requires validating that MCS implementations are sound with respect to their specifications. MCS testing techniques already exist and have found important bugs, but for modern cross-platform GPU frameworks they remain insufficient in three areas: measuring the effectiveness of MCS testing, scaling MCS testing to reach the diversity of GPUs available today, and analyzing the intersection of MCS and memory safety guarantees to increase the security of modern GPU frameworks.

\subsection{Effectiveness of MCS Testing}

For a programmer to be able to reason about the outcomes of their program, they assume that the platform they are running on, i.e., the combination of compilers, drivers, and hardware, implements the MCS soundly. Unfortunately, this is not always the case; MCS testing strategies have found numerous bugs in runtime platforms~\cite{kirkham2020foundations,alglave2015gpu,manerkar2016counterexamples,manerkar2017rtlcheck,alglave2010fences}. Early MCS testing work executed sequences of memory accesses~\cite{archtest,hangal2004tsotool} and then analyzed a trace of the observed values.  More recent work uses formal specifications to derive small concurrent tests like the message passing program in \myfig\ref{fig:intro-mp}. These programs, called \textit{litmus tests}, are then executed for many iterations, checking for a weak behavior at each iteration~\cite{alglave2011litmus}.

The effectiveness of litmus testing approaches is \textit{extremely} sensitive to the testing environment, i.e.\ the context around the litmus test that creates stress on the platform. When testing, it is impossible to tell if an unobserved illegal execution, e.g., an execution of \myfig\ref{fig:intro-mprelacq} where the weak behavior is observed, is impossible or if it is simply rare and was not exposed by the tests. Thus, effective MCS testing requires \textit{tuning} testing environments, i.e., adapting the parameters of the environment to be more effective. However, legitimate MCS violations are exceedingly rare, and thus cannot be used as a metric for such tuning efforts. Existing approaches have tuned their testing environments by maximizing the number of weak, but allowed, behaviors of some litmus tests, yet have not provided justification on why this metric is valid. Existing MCS testing techniques also do not provide an empirical confidence metric in their ability to find MCS bugs. Consequently, it is difficult to gauge whether it is worth the time to perform more testing, especially when running in time-constrained environments such as a conformance test suite. The result of these limitations has been a lack of large, rigorously backed test suites for validating MCS implementations, with existing conformance test suites like those for the Vulkan~\cite{khronos_vk_gl_cts} and WebGPU~\cite{gpuweb_cts} containing few to no ad hoc MCS tests running in unvalidated testing environments.

\subsection{Scale of MCS Testing}

Previous GPU testing work had limited scope, testing only a small number of devices~\cite{kirkham2020foundations}, with the largest study testing eight devices~\cite{alglave2015gpu}. These approaches did not scale due to the difficulty of portable GPU application development and deployment, e.g., while frameworks like OpenCL~\cite{opencl} are portable in theory, there are many difficulties in practice~\cite{sorensen2016hitchhiker}.  Consequently, little is known about the MCS conformance and weak behavior profiles \emph{at large}.  This is especially problematic as portable GPU frameworks depend upon many layers and environments (e.g., architectures, compilers, runtimes, operating systems, etc.); it is difficult to extrapolate insights from a small number of platforms tested in controlled environments to the diverse universe of deployed GPUs.

Apart from conformance, a device's weak behavior \textit{profile}, i.e., the frequency at which allowed weak behaviors occur and how system stress influences this frequency, is also a useful metric. For example, this can be useful in developing conformance testing strategies and enables developers to reason about tradeoffs between accuracy and performance in approximate computing that judiciously elides synchronization~\cite{niu2011hogwild,renganarayana2012programming,samadi2013sage}. Focused collection of GPU weak behavior profiles also enables more rigorous analysis of GPU architecture, helps determine how precisely language-level MCSs describe deployed implementations, and can be used as the starting point to reveal new hardware side-channels.

\subsection{Intersection of the MCS and Memory Safety}

Over the past several decades, CPU-based languages have been designed to provide \textit{memory safety}, and runtimes, e.g., operating systems, have developed protections for memory isolation between processes, e.g., via segfaults. Some of the most notorious exploits and vulnerabilities found in recent decades were related to memory safety~\cite{heartbleed,trident,slammer-worm}, and large-scale corporate studies have found that the majority of all vulnerabilities are related to memory safety~\cite{microsoft2019safer,google2019queue}.

Due to their lean, performance oriented implementations and relatively more recent adoption, barely any GPU runtimes have these defensive measures in place. One area where memory safety is extremely important is in the browser, where web pages can execute code from untrusted sources which may maliciously try to access data from other tabs or applications on a user's machine. New GPU frameworks like WebGPU provide low-level access to the GPU from the browser but must also ensure that GPU code running in one tab cannot access memory from another tab. This problem is complicated by the fact that the WebGPU implementation does not directly control the final machine-code that runs on the GPU; rather, it goes through several intermediate languages and frameworks each with its own specification and implementations, creating a fundamental mismatch between the WGSL source language and lower-level languages.

While WGSL may promise memory safety for untrusted code, lower-level languages may have weaker or undefined semantics for racy memory accesses. In ordinary unsafe GPU programming models, where an application must be explicitly installed and likely has a more trustworthy provenance, avoiding races can be left to the programmer, but in a browser programming model, data races cannot be allowed to invalidate safety-critical program slices such as bounds checks, index computations, or buffer isolation logic. Thus, preserving memory safety requires reasoning not only about explicit guardrails such as clamping out-of-bounds accesses, but also about whether the MCS gives those guardrails stable meaning after compilation. While previous work on CPU-based languages has sought to constrain the effects of data races~\cite{dolan2018bounding} and develop safe languages that run in the browser~\cite{watt2020repairing,watt2019weakening}, no research has focused on the complexity of maintaining GPU memory safety in the presence of data races.

\section{Dissertation Overview}
\label{sec:intro-contrib}

This thesis directly addresses the need for more principled and scalable techniques to validate the correctness and safety of MCS implementations on GPUs. First, \mychaplong\ref{chap:background} provides background on GPU programming and memory consistency. Then, each of the next three chapters describes the core original contributions of this thesis, which are introduced briefly here.

\paragraph{MC Mutants.}  \mychaplong\ref{chap:mc-mutants} introduces \textit{MC Mutants}, a methodology for evaluating the effectiveness of MCS testing, and evaluates it on several GPUs. MC Mutants is a mutation testing approach for evaluating MCS testing environments, mutating MCS tests such that the mutants simulate bugs that might occur. A testing environment can then be evaluated using a mutation score. MC Mutants is used to evaluate a novel parallel testing environment, implemented in WebGPU, which improves testing speed compared to previous approaches by three order of magnitude. MC Mutants also provides an MCS testing confidence strategy that is parameterized over a time budget and confidence threshold, showing that GPU MCS tests using the parallel testing environment need to run for only $\sim$1 minute per test, compared to hours or even days by previous methods. MC Mutants is evaluated on four GPUs from four vendors, identifying two MCS bugs in WebGPU implementations and leading to the official WebGPU conformance test suite adopting tests developed using MC Mutants.

\paragraph{GPUHarbor.} \mychaplong\ref{chap:gpuharbor} presents \textit{GPUHarbor}, a wide-scale GPU MCS testing tool with a web interface and an Android app. GPUHarbor is used to run a testing campaign that checks conformance to an MCS and characterizes weak behaviors. GPUHarbor was advertised on forums and social media, leading to data collection from 106 devices spanning 7 vendors. In terms of devices tested, this constitutes the largest study on weak memory behaviors by at least $10\times$, and moreover the conformance tests run during the testing campaign identified two new bugs on embedded Arm and NVIDIA devices. Analyzing characterization data yields many insights, including quantifying and comparing weak behavior occurrence rates (e.g., AMD GPUs show $25.3\times$ more weak behaviors on average than Intel).
The results from the GPUHarbor study are also discussed in relation to the impact they have on software development for GPUs.

The techniques used to scale GPU MCS testing are applied in several other contributions as well. Focused testing of modern NVIDIA GPUs is done to validate their memory consistency implementations and refine the understanding of what MCS properties they can support. Parallel MCS testing environments also show they are capable of revealing hardware side-channels in various GPUs, as stress from one process is able to trigger weak behaviors in a separate process. A low-capacity side-channel is developed as a proof-of-concept, while these initial observations also form the foundation for deeper analysis into the security guarantees provided by NVIDIA's Multi-Instance-GPU (MIG) technology~\cite{nvidia2020ampere}.

\paragraph{SafeRace.} \mychaplong\ref{chap:saferace} introduces \textit{SafeRace}, consisting of a set of tools that are used to perform targeted security assessments of existing systems and several proposals to increase the security of WebGPU's programming language, WGSL. SafeRace is motivated by the identification of a \textit{specification vulnerability} in WebGPU, where code with data races can be compiled to intermediate representations in which an optimizing compiler could legitimately remove memory safety guardrails. While the threat assessment shows that this vulnerability does not appear to be exploitable on current systems, it creates a ``ticking time bomb'', especially as compilers in this area are rapidly evolving. Given this, the SafeRace Memory Safety Guarantee (SMSG) is introduced, two components that preserve memory safety in the WGSL execution stack even in the presence of data races. The first component specifies that program slices contributing to memory indexing must be race free and is implemented via a compiler pass for WGSL programs. The second component is a requirement on intermediate representations that limits the effects of data races so that they cannot impact race-free program slices.

The first component is evaluated by implementing a new transformation in a WGSL compiler and showing that it has minimal performance overhead on existing WebGPU applications. The second component is evaluated by performing a fuzzing campaign of 81 hours across 21 compilation stacks that shows lower-level GPU frameworks could relatively straightforwardly adopt the necessary requirements. The threat assessments discover GPU memory isolation vulnerabilities in Apple and AMD GPUs, as well as a security-critical miscompilation of WGSL in a pre-release version of Firefox.


Finally, \mychaplong\ref{chap:conclusion} summarizes all findings and provides a discussion of future work.


\section{Publications}
\label{sec:intro-pub}

The three major contributions in this thesis, MC Mutants, GPUHarbor, and SafeRace, represent research projects that I led and have been published in peer-reviewed venues. These works have also led to several related and follow-on works, which have either been published in peer-reviewed venues, submitted as preprints to open-access research repositories, or are currently under submission. Organized by the nature of my contribution, the full list of publications is:

\paragraph{Lead Author}
\begin{itemize}
  \item \textbf{Reese Levine}, Tianhao Guo, Mingun Cho, Alan Baker, Raph Levien, David Neto, Andi Quinn, and Tyler Sorensen. 2023. \textit{MC Mutants: Evaluating and Improving Testing for Memory Consistency Specifications}. In Proceedings of the 28th ACM International Conference on Architectural Support for Programming Languages and Operating Systems, Volume 2 (ASPLOS 2023)~\cite{levine2023mcmutants}.
  \begin{itemize}
    \item This work received both \textbf{Distinguished Paper} and \textbf{Distinguished Artifact} awards.
    \item The contributions in this work are presented in \mychaplong\ref{chap:mc-mutants}.
  \end{itemize}

  \item \textbf{Reese Levine}, Mingun Cho, Devon McKee, Andi Quinn, and Tyler Sorensen. 2023. \textit{GPUHarbor: Testing GPU Memory Consistency at Large (Experience Paper)}. In Proceedings of the 32nd ACM SIGSOFT International Symposium on Software Testing and Analysis (ISSTA 2023)~\cite{levine2023gpuharbor}.
  \begin{itemize}
    \item This work received a \textbf{Distinguished Artifact} award.
    \item The contributions in this work are presented in \mychaplong\ref{chap:gpuharbor}.
  \end{itemize}

  \item \textbf{Reese Levine}, Ashley Lee, Neha Abbas, Kyle Little, and Tyler Sorensen. 2025. \textit{SafeRace: Assessing and Addressing WebGPU Memory Safety in the Presence of Data Races}. Proc. ACM Program. Lang. 9, OOPSLA2, Article 297 (October 2025)~\cite{levine2025saferace}.
  \begin{itemize}
    \item The contributions in this work are presented in \mychaplong\ref{chap:saferace}.
  \end{itemize}

  \item \textbf{Reese Levine}, Rithik Sharma, Nikhil Jain, Abhijit Ramesh, Zheyuan Chen, Neha Abbas, James Contini, and Tyler Sorensen. 2026. \textit{Llamas on the Web: Memory-Efficient, Performance-Portable, and Multi-Precision LLM Inference with WebGPU}. arXiv preprint arXiv:2605.20706~\cite{levine2026llamaweb}.
  \begin{itemize}
    \item This work is not directly discussed in this thesis, but follows from the research on increasing the reliability and security of WebGPU, and is briefly mentioned in \mysec\ref{sec:conclusion-future-work}.
  \end{itemize}
\end{itemize}

\paragraph{Co-Author}
\begin{itemize}
  \item Sean Siddens, Sanya Srivastava, \textbf{Reese Levine},\\
  Josiah Dykstra, and Tyler Sorensen. 2026.\\
  \textit{Memory DisOrder: Memory Re-orderings as a Timerless Side-channel}. arXiv preprint arXiv:2601.08770~\cite{siddens2026disorder}.
  \begin{itemize}
    \item My contributions to this work are discussed in \mysec\ref{sec:gpuharbor-sidechannel}.
  \end{itemize}

  \item Cheng Gu, \textbf{Reese Levine}, Zhenkai Zhang, Tyler Sorensen, and Yanan Guo. 2026. \textit{Behind Bars: A Side-Channel Attack on NVIDIA MIG Cache Partitioning Using Memory Barriers}. To appear in the 35th USENIX Security Symposium~\cite{gu2026behindbars}.
  \begin{itemize}
    \item My contributions to this work are discussed in \mysec\ref{sec:gpuharbor-sidechannel}.
  \end{itemize}

  \item Soham Bagchi, Sanya Srivastava, \textbf{Reese Levine}, Tyler Sorensen, Ryan Stutsman, and Vijay Nagarajan. 2026. \textit{Consistency and Coherence of the NVIDIA Grace-Hopper Superchip}. In ACM SIGPLAN International Symposium on Memory Management (ISMM 2026)~\cite{bagchi2026gh}.
  \begin{itemize}
    \item My contributions to this work are discussed in \mysec\ref{sec:gpuharbor-nvidia-grace-hopper}.
  \end{itemize}

  \item Dennis Sprokholt, \textbf{Reese Levine}, Susmit Sarkar, Andr\'es Goens, Vijay Nagarajan, and Soham Chakraborty. \textit{Beyond the Fence: Sound Compilation for GPU and Heterogeneous CPU+GPU Systems}. Under submission.
  \begin{itemize}
    \item My contributions to this work are discussed in \mysec\ref{sec:gpuharbor-nvidia-beyond-fence}.
  \end{itemize}

\end{itemize}

\paragraph{Workshop and Article}
\begin{itemize}
  \item \textbf{Reese Levine} and Tyler Sorensen. 2023. \textit{Probabilistic Memory Consistency Specifications}. In ACM Young Architect Workshop~\cite{levine2023pmcs}.
  \begin{itemize}
    \item The contributions in this work are discussed in \mysec\ref{sec:conclusion-future-work}.
  \end{itemize}

  \item Arkaprava Basu, Natalia Gavrilenko, Keijo Heljanko, \textbf{Reese Levine}, Ajay Ashok Nayak, Hern\'an Ponce de Le\'on, Tyler Sorensen, and Haining Tong. 2025. \textit{GPU Memory Consistency: Specifications, Testing, and Opportunities for Performance Tooling}. ACM SIGARCH Blog~\cite{basu2025gpu-mcs-blog}.
  \begin{itemize}
    \item My contributions to this work are discussed in \mysec\ref{sec:gpuharbor-broader-impact}.
  \end{itemize}

\end{itemize}

\paragraph{Funding Acknowledgement.} I have been fortunate to have been supported in part by a National Defense Science and Engineering Graduate (NDSEG) Fellowship.
